\section{Which transition energies preserve the reduction?}
\label{sec:transitionenergy}
Zero transition energy ensures that total service consumption is independent
of the matching. A controlled extension permits a common energy cost per
matched pair. This is a compulsory service-transition task, performed over
a common positive-duration, noncharging interval after the early window
and before the B batch. It is an abstract extension of the common-window
model, not a claim about arbitrary road movements or charging access.

\begin{corollary}[Common transition energy]\label{cor:commontransition}
Retain the assumptions of Theorem~\ref{thm:lossyfrontier}, except that each
matched A--B pair consumes the same additional energy $d\ge0$ during that
interval, with nonnegative consumption throughout. Insert this interval
after early exchange and translate the B batch and terminal window later
without changing their durations or budgets: for example, A on $[0,1]$,
early exchange on $[1,2]$, transition on $[2,3]$, B on $[3,4]$ and terminal
exchange on $[4,5]$. Singletons perform no
transition task. Restrict the allowed graph to edges with $b_j+d\le B$.
Write $S=\sum_i a_i+\sum_j b_j$, $P=\sum_i a_i$,
$S_m=S+md$, $N_m=2n-m$, and
\[
 r_m=\min_{|M|=m}\sum_{(i,j)\in M}(a_i+b_j+d-B)_+,
 \quad U=\min\{P/\alpha,K\},
 \quad \kappa=\frac1{\alpha\delta}-1.
\]
For $m=0,\ldots,\nu$ in this restricted graph, the efficient curve is
\[
 L_m(E)=\frac{S_m}{\alpha}-E+
       \kappa\left(\frac{r_m}{\alpha}-E\right)_+,
 \qquad t_m\le E\le U,
\]
where
\begin{align*}
 t_m=\max\Biggl\{&
 (\alpha^{-1}-\delta)r_m+\delta(S_m-N_mB),\quad
 2r_m/\alpha-K,\\[-2pt]
 &S_m/\alpha-H,\quad
 \frac{S_m/\alpha+\kappa r_m/\alpha-H}{1+\kappa}
 \Biggr\}.
\end{align*}
Intervals with $t_m>U$ are infeasible. There are at most two efficient
segments per cardinality. Under coordinatewise nondecreasing supply costs,
these curves suffice for physical and convexified planning values and for
nonnegative posted-price value queries, as in
Section~\ref{sec:lossyalgorithm}. They do not enumerate every response tie
or the full physical projection. The sequence $t_m$ is discrete convex.
\end{corollary}

The proof uses the same interval allocation as
Lemma~\ref{lem:lossyfixed}: a paired bus now has early battery-exchange box
$[a_i+b_j+d-B,a_i]$. Its lower endpoints sum to $S_m-N_mB$ and its positive
lower endpoints to the matching requirement. Thus both the requirement and
the terminal energy account include the task. Appendix~\ref{app:transitionenergy}
gives the details, including capped paid-idle dominance. No lossless scalar
support criterion is transferred by this corollary. The separate full-polygon
qualification includes one common-task literal and checks the efficient lower
threshold as well; the continuous-domain claim rests on the proof.

\begin{corollary}[Complete polygons with a common transition task]
\label{cor:commontransitionfull}
Retain the timing, restricted graph $b_j+d\le B$ and definitions
$S_m,N_m,r_m,P$ of Corollary~\ref{cor:commontransition}, and permit
the finite early actions of Theorem~\ref{thm:fulllossyprojection}.
For service-assigned fleets, the complete physical projection is the union
over $m=0,\ldots,\nu$ of
\[
 \left\{\left(C-X,\frac{S_m-Q}{\alpha},fN_m\right):
          (C,X)\in W_m^{d}\right\},
\]
where
\[
\begin{split}
 W_m^{d}=\bigl\{(C,X):\;&C\ge r_m/\alpha,\quad X\ge0,
                  \quad C+X\le K,\\
 &\max\{S_m-N_mB,S_m-\alpha H\}
       \le Q=\alpha C-X/\delta\le P\bigr\}.
\end{split}
\]
Each nonempty polygon has at most five vertices. Their affine images
generate the complete-plan hull. Keeping these full polygons preserves
arbitrary projected-objective values, including negative-price queries;
no monotonicity assumption is required.
\end{corollary}
\begin{proof}
The paired exchange box is $[a_i+b_j+d-B,a_i]$; singleton boxes are
unchanged. Their lower endpoints sum to $S_m-N_mB$, upper endpoints to
$P$, and positive lower endpoints to the matching requirement.
The finite-cycle construction of Theorem~\ref{thm:fulllossyprojection}
occurs before the transition task and returns to the starting early SOC.
Nonnegative task and B-service consumption then preserves intermediate
SOC feasibility. At fixed $m$, $S_m$ is independent of the matching,
so a matching attaining $r_m$ contains every other matching's polygon.
The full-set and hull conclusions follow.
\end{proof}
For $d>0$ this is not a literal paid-idle elimination result. Deleting
an edge while replacing its paired bus and an idle bus by two singletons
reduces total consumption by $d$. It changes terminal load even at the
same $C,X$, so only the earlier monotone-value dominance claim follows.
The statement also excludes edge-dependent task energies, charging during
the task and altered access policies. One common-task literal passed the
separate full-polygon qualification; it does not qualify arbitrary task models.

\paragraph{Varying edge energies need more information.}
The failure already occurs with charging only; disable early exports in
this example. Consider $n=2$, $a=b=(13,1)$, $B=20$, unit efficiency, early budget $K=14$
and terminal budget $H=40$. Every edge is allowed, with transition energies
$\left(\begin{smallmatrix}0&6\\6&0\end{smallmatrix}\right)$.
Use A services on $[0,1]$, early charging on $[1,2]$, positive-duration
noncharging transitions on $[2,3]$, B services on $[3,4]$ and terminal
charging on $[4,5]$. One connector of power 14 early and 40 terminal,
with adequate acceptance, serially realizes the stated budgets. Access and
other nonservice energy are zero. Transition tasks are compulsory; no
metric-road realization or substitute route is assumed.
\begin{center}
\begin{tabular}{lccc}
\toprule
Full matching & Minimum early requirement & Total energy & Load segment\\
\midrule
Identity & $6$ & $28$ & $(E,28-E),\ 6\le E\le14$\\
Cross & $0$ & $40$ & $(E,40-E),\ 0\le E\le14$\\
\bottomrule
\end{tabular}
\end{center}
At $E=6$ the identity's terminal load 22 is smaller than every cross load
with early purchase at most 6. Selecting only the minimum-requirement
matching therefore loses an efficient branch. For the fixed monotone
convex supply cost $F(E,L)=E+L$, its two-bus value is $2f+28$ instead of
the cross value $2f+40$. Allowing singleton duties does not repair the
error when $f>0$: the true value is $2f+28$, whereas even the best
minimum-requirement choice at every cardinality gives
\[
 \min\{2f+40,3f+28,4f+28\}
 =2f+28+\min\{12,f\}>2f+28.
\]
This elementary counterexample isolates the lost assignment-independence
of total energy. It is neither a general hardness result nor a public-case
or empirical finding.
Appendix~\ref{app:exponentialtransition} extends the construction to
$2^k$ efficient physical pieces at a single fleet size, each necessary for
some monotone convex planning objective. Its numerical magnitudes grow
with $k$, and the hull remains simple; it proves neither hardness nor an
exponential requirement for every symbolic representation.
For one smooth monotone convex objective in that family, the exact gap is
$9/(2K)>0$ even with all these alternatives available. The gap decreases
with the family's charging budget $K$; there is no uniform-size claim.
Appendix~\ref{app:completegraphtransition} then adds every interblock edge
and admits every perfect matching, while preserving that gap. Thus graph
disconnection is unnecessary for this example. The compulsory-task model
and fleet cap remain essential stated restrictions.

% Invoke once in the appendix, after this file has been input.
\newcommand{\TransitionEnergyProofs}{%
\section{Proof of the common-transition extension}\label{app:transitionenergy}
\begin{proof}[Proof of Corollary~\ref{cor:commontransition}]
For a fixed matching, paired, A-only and B-only early battery-exchange
boxes are, respectively,
\[
 [a_i+b_j+d-B,a_i],\qquad [a_i-B,a_i],\qquad[b_j-B,0].
\]
The edge condition makes the first box nonempty. Capacity after A gives
the upper endpoints; nonnegative SOC after the transition and B gives the
lower endpoints. Nonnegative consumption during the post-early portion
then guarantees all its intermediate SOC constraints. The upper endpoints
sum to $P$, all lower endpoints to $T=S_m-N_mB$, and their positive parts
to $R=\sum_{(i,j)\in M}(a_i+b_j+d-B)_+$.

The allocation argument of Lemma~\ref{lem:lossyfixed} therefore gives
minimum export $D_R(E)=(R/\alpha-E)_+$ and terminal load
$S_m/\alpha-E+\kappa D_R(E)$. The available negative battery exchange,
early gross budget and terminal budget give the four lower bounds in
$t_m$ with $R$ in place of $r_m$; early capacity and gross budget give
$E\le\min(P/\alpha,K)$. Every point in that interval is attained by
distributing positive or negative exchange over the boxes, executing early
exchanges serially, performing the noncharging task and B service, and
refilling serially at the terminal. Extra cycles weakly increase terminal
load and resource use, so monotone value objectives permit their removal.

At fixed $m$, total energy $S_m$ and fleet $N_m$ are constant across
matchings. Both the terminal curve and each threshold term are
nondecreasing in $R$, since $\alpha^{-1}-\delta\ge0$ and $\kappa\ge0$.
A matching attaining $r_m$ consequently dominates all other size-$m$
curves. Componentwise replacement also proves the convexified-value
claim. Classical cardinality-matching convexity makes $r_m$ discrete
convex, while $S_m$ and $S_m-N_mB$ are affine. Each threshold term is
discrete convex and so is their maximum: extend each term by its convex
linear interpolant, take the convex maximum, and restrict to integers.

For completeness, suppose $k$ paid idle buses are admitted with total
$N=2n-m+k\le2n$, all available early. Their boxes are $[-B,0]$.
Delete any $k\le m$ matching edges and replace each deleted pair plus
one idle bus by its two singleton services. Fleet and fee stay unchanged,
$P$ stays unchanged, total energy drops by $kd$, and $R$ weakly decreases.
Every threshold term weakly decreases and the terminal curve weakly
decreases at the same $E$. The no-idle representative therefore dominates
under monotone economics. This proves efficient-value dominance, not
equality of feasible load sets; it makes no claim about uncapped idle fleets.
\end{proof}
}
